% !TEX program = pdflatex
\documentclass[10pt]{article}

\usepackage[utf8]{inputenc}
\usepackage[T1]{fontenc}
\usepackage[english]{babel}
\usepackage{lmodern}
\usepackage{microtype}
\usepackage[margin=0.72in]{geometry}
\usepackage{amsmath,amssymb,amsthm,mathtools}
\usepackage{booktabs,array}
\usepackage{xcolor}
\usepackage{enumitem}
\usepackage[most]{tcolorbox}
\usepackage{fancyhdr}
\usepackage[hidelinks]{hyperref}

\definecolor{CourseBlue}{HTML}{1A3C68}
\definecolor{CourseGreen}{HTML}{216E39}
\definecolor{CourseGold}{HTML}{B8860B}
\definecolor{CourseRed}{HTML}{9A2C2C}
\definecolor{SoftBlue}{HTML}{F2F6FA}
\definecolor{SoftGreen}{HTML}{F1F8F3}
\definecolor{SoftGold}{HTML}{FFF8E8}
\definecolor{SoftRed}{HTML}{FFF3F3}

\setlength{\parindent}{0pt}
\setlength{\parskip}{0.32em}
\setlist[itemize]{leftmargin=1.5em,itemsep=0.12em,topsep=0.2em}
\setlist[enumerate]{leftmargin=1.7em,itemsep=0.12em,topsep=0.2em}
\allowdisplaybreaks

\pagestyle{fancy}
\fancyhf{}
\lhead{Stochastic Calculus and Option Pricing}
\rhead{Lecture 7 Notes}
\cfoot{\thepage}
\renewcommand{\headrulewidth}{0.3pt}

\newcommand{\E}{\mathbb E}
\newcommand{\Pp}{\mathbb P}
\newcommand{\Q}{\mathbb Q}
\newcommand{\R}{\mathbb R}
\newcommand{\1}{\mathbf 1}
\newcommand{\cF}{\mathcal F}
\newcommand{\dd}{\,\mathrm d}
\newcommand{\Nlaw}{\mathcal N}
\newcommand{\esssup}{\operatorname*{ess\,sup}}
\DeclareMathOperator{\Var}{Var}
\DeclareMathOperator{\Cov}{Cov}
\DeclareMathOperator{\rank}{rank}

\theoremstyle{definition}
\newtheorem{definition}{Definition}[section]
\newtheorem{proposition}[definition]{Proposition}
\newtheorem{theorem}[definition]{Theorem}
\newtheorem{lemma}[definition]{Lemma}
\newtheorem{corollary}[definition]{Corollary}
\newtheorem{remark}[definition]{Remark}

\newtcolorbox{boardexample}[1]{
  enhanced,breakable,colback=SoftBlue,colframe=CourseBlue,
  boxrule=0.55pt,arc=1mm,left=1.5mm,right=1.5mm,top=1mm,bottom=1mm,
  before skip=0.55em,after skip=0.55em,
  title={Example --- #1},fonttitle=\bfseries}
\newtcolorbox{studentexercise}[1]{
  enhanced,breakable,colback=SoftGold,colframe=CourseGold,
  boxrule=0.55pt,arc=1mm,left=1.5mm,right=1.5mm,top=1mm,bottom=1mm,
  before skip=0.55em,after skip=0.35em,
  title={Exercise --- #1},fonttitle=\bfseries}
\newtcolorbox{shortsolution}[1][]{
  enhanced,breakable,colback=SoftGreen,colframe=CourseGreen,
  boxrule=0.55pt,arc=1mm,left=1.5mm,right=1.5mm,top=1mm,bottom=1mm,
  before skip=0.25em,after skip=0.6em,
  title={Short solution#1},fonttitle=\bfseries}
\newtcolorbox{keyidea}[1]{
  enhanced,breakable,colback=white,colframe=CourseGreen,
  boxrule=0.7pt,arc=1mm,left=1.5mm,right=1.5mm,top=1mm,bottom=1mm,
  before skip=0.5em,after skip=0.5em,
  title={#1},fonttitle=\bfseries}
\newtcolorbox{warningbox}[1]{
  enhanced,breakable,colback=SoftRed,colframe=CourseRed,
  boxrule=0.65pt,arc=1mm,left=1.5mm,right=1.5mm,top=1mm,bottom=1mm,
  before skip=0.5em,after skip=0.5em,
  title={#1},fonttitle=\bfseries}
\newtcolorbox{assumptionbox}[1]{
  enhanced,breakable,colback=SoftGold,colframe=CourseGold,
  boxrule=0.65pt,arc=1mm,left=1.5mm,right=1.5mm,top=1mm,bottom=1mm,
  before skip=0.5em,after skip=0.5em,
  title={#1},fonttitle=\bfseries}

\title{\textbf{Martingale Pricing, Hedging, and No-Arbitrage}\\
\large Lecture 7 Notes --- From Martingale Prices to Replication and Market Completeness}
\author{}
\date{}

\begin{document}
\maketitle
\vspace{-1.4em}

\begin{keyidea}{Goals and logical structure}
Lecture 6 constructed the pricing measure. Lecture 7 follows one chain:
\[
\begin{aligned}
\text{risk-neutral stock dynamics}
&\longrightarrow \text{discounted self-financing gains and }\Q\text{-prices}\\
&\longrightarrow \text{dividend-paying call replication under }\Pp\\
&\longrightarrow \text{dividend-adjusted pricing and a new numeraire}\\
&\longrightarrow \text{martingale hedges and two-risk spanning}\\
&\longrightarrow \text{completeness and path dependence}\\
&\longrightarrow \text{It\^o, Girsanov, reflection, and stopping-time revision}.
\end{aligned}
\]
The recurring discipline is to separate four questions: which numeraire is
used, under which measure the discounted gains are martingales, whether the
claim is replicable, and whether the relevant local martingales are true
martingales.
\end{keyidea}

% ============================================================
\section{The Risk-Neutral Starting Point}
% ============================================================

\subsection{The five-line bridge from Lecture 6}

For the constant-coefficient stock model under \(\Pp\), let
\(\lambda=(\mu-r)/\sigma\). The preceding lecture ended with the chain
\[
Z_t=e^{-\lambda W_t^{\Pp}-\lambda^2t/2},\qquad
\frac{\dd\Q}{\dd\Pp}\bigg|_{\cF_T}=Z_T,\qquad
W_t^{\Q}=W_t^{\Pp}+\lambda t,
\]
\[
\frac{\dd S_t}{S_t}=r\dd t+\sigma\dd W_t^{\Q},\qquad
\dd(e^{-rt}S_t)=\sigma e^{-rt}S_t\dd W_t^{\Q}.
\]
The density changes probabilities; the last expression identifies a
stochastic integral. Its true-martingale status needs a separate
integrability argument.

For now, use the money-market account \(B^{(0)}\) as the numeraire. Under
its equivalent martingale measure, discounted traded prices or total
gains are \(\Q\)-local martingales; under the conditions used for
valuation, they are true martingales. We change numeraire after treating
dividends.

\subsection{Why the discounted stock is a local martingale}

For a non-dividend-paying stock under the risk-neutral measure, suppose
\[
\dd S_t=r_tS_t\dd t+\sigma_tS_t\dd W_t^\Q,
\qquad
D_t=\exp\left(-\int_0^t r_u\dd u\right).
\]
The product rule removes the drift:
\[
\boxed{
D_tS_t=S_0+\int_0^tD_u\sigma_uS_u\dd W_u^\Q.}
\]
A well-defined stochastic integral with respect to Brownian motion is a
local martingale. This statement by itself does not assert that the
discounted price is a true martingale; that requires an additional
integrability argument.

The stock is now a local martingale after discounting. We next show how
self-financing gains lead to a conditional-expectation price whenever a
portfolio replicates a payoff.

% ============================================================
\section{Self-Financing Leads to Martingale Prices}
% ============================================================

We first derive the martingale valuation identity for any replicating
self-financing strategy. We will then construct such a strategy explicitly
for a call on a dividend-paying stock.

\subsection{Discounted self-financing gains}

Let the money-market account be
\[
B_t^{(0)}=\exp\left(\int_0^t r_u\dd u\right),
\qquad
D_t=\frac{1}{B_t^{(0)}}.
\]
For this derivation, suppose the risky assets pay no dividends. Consider
\(S^1,\ldots,S^m\) and a predictable strategy
\(\phi=(\phi^0,\phi^1,\ldots,\phi^m)\) with wealth
\[
X_t=\phi_t^0B_t^{(0)}+\sum_{i=1}^m\phi_t^iS_t^i.
\]
The strategy is self-financing when
\[
\dd X_t=\phi_t^0\dd B_t^{(0)}
        +\sum_{i=1}^m\phi_t^i\dd S_t^i.
\]
Define discounted prices
\[
\bar X_t=D_tX_t,
\qquad
\bar S_t^i=D_tS_t^i.
\]

\begin{proposition}[Discounted gain identity]
For a self-financing strategy,
\[
\boxed{\dd\bar X_t=\sum_{i=1}^m\phi_t^i\dd\bar S_t^i.}
\]
Equivalently,
\[
\bar X_t=\bar X_0+\sum_{i=1}^m
\int_0^t\phi_u^i\dd\bar S_u^i.
\]
\end{proposition}

\begin{proof}
Since \(D\) has finite variation,
\[
\dd(D_tX_t)=D_t\dd X_t+X_t\dd D_t,
\qquad
\dd D_t=-r_tD_t\dd t.
\]
Use \(\dd B_t^{(0)}=r_tB_t^{(0)}\dd t\) and
\(\phi_t^0B_t^{(0)}=X_t-\sum_i\phi_t^iS_t^i\). The financing terms cancel,
leaving precisely the discounted risky-asset gains.
\end{proof}

\begin{studentexercise}{recover the cancellation in one dimension}
For \(X_t=\Delta_tS_t+\beta_tB_t^{(0)}\), start from
\(\dd X_t=\Delta_t\dd S_t+\beta_t\dd B_t^{(0)}\) and derive
\(\dd(D_tX_t)=\Delta_t\dd(D_tS_t)\).
\end{studentexercise}
\begin{shortsolution}
\[
\begin{aligned}
\dd(DX)
&=D\Delta\dd S+D\beta rB^{(0)}\dd t-rDX\dd t\\
&=\Delta(D\dd S-rDS\dd t)
=\Delta\dd(DS).
\end{aligned}
\]
The equality uses \(X=\Delta S+\beta B^{(0)}\). Discounting removes the
cash account from the gain equation, not from the portfolio value.
\end{shortsolution}

\subsection{The martingale valuation identity}

Fix an equivalent martingale measure \(\Q\) associated with the money-market
numeraire. By definition, discounted traded gains are \(\Q\)-local
martingales. Under the required admissibility and integrability assumptions,
the discounted wealth of a self-financing strategy is a true martingale.

\begin{theorem}[Pricing a replicable claim]
Let \(H\) be \(\cF_T\)-measurable and suppose a self-financing strategy
replicates it:
\[
X_T=H.
\]
If \(\bar X=DX\) is a true \(\Q\)-martingale, then its value is
\[
\boxed{
V_t=X_t
=B_t^{(0)}\E^\Q\left[\frac{H}{B_T^{(0)}}\middle|\cF_t\right]
=\frac{1}{D_t}\E^\Q[D_TH\mid\cF_t].}
\]
When rates are deterministic,
\[
\boxed{
V_t=\E^\Q\left[
\exp\left(-\int_t^T r_u\dd u\right)H
\middle|\cF_t\right].}
\]
\end{theorem}

\begin{proof}
The discounted gain identity makes \(D_tX_t\) a stochastic integral with
respect to discounted traded martingales. Once it is a true martingale,
\[
D_tX_t=\E^\Q[D_TX_T\mid\cF_t]
=\E^\Q[D_TH\mid\cF_t].
\]
Multiply by \(D_t^{-1}=B_t^{(0)}\).
\end{proof}

\begin{warningbox}{Local martingale is not enough}
No drift is an algebraic statement; pricing by conditional expectation is an
integrability statement. A discounted self-financing wealth process is first
a local martingale. Admissibility, square integrability, nonnegativity, or a
problem-specific argument is needed before asserting
\(\bar X_t=\E^\Q[\bar X_T\mid\cF_t]\).
\end{warningbox}

\begin{boardexample}{a unit payoff and a terminal share}
Assume constant \(r\) and no dividends. For \(H=1\),
\[
V_t=e^{-r(T-t)}.
\]
For \(H=S_T\),
\[
V_t=B_t^{(0)}
\E^\Q\left[\frac{S_T}{B_T^{(0)}}\middle|\cF_t\right]
=S_t,
\]
because \(S/B^{(0)}\) is a \(\Q\)-martingale. A payoff of one bank-account
unit at \(T\) would instead cost \(B_t^{(0)}\). These values can also be
read directly from their replicating portfolios.
\end{boardexample}



\subsection{The lognormal law, digitals, and stock-weighted tails}

For a replicable integrable payoff \(H\), when discounted wealth is a
true martingale, constant rates give
\[
\boxed{V_t=e^{-r(T-t)}\E^\Q[H\mid\cF_t].}
\]
Under \(\Q\), conditional on \(\cF_t\), the Brownian increment
\(W_T^\Q-W_t^\Q\) is normal with mean zero and variance
\(\tau=T-t\). Hence
\[
S_T=S_t\exp\!\left((r-\tfrac12\sigma^2)\tau
+\sigma(W_T^\Q-W_t^\Q)\right).
\]
Define
\[
d_2=\frac{\log(S_t/K)+(r-\sigma^2/2)\tau}{\sigma\sqrt\tau},
\qquad d_1=d_2+\sigma\sqrt\tau.
\]
A cash-or-nothing digital paying
\(\1_{\{S_T>K\}}\) has value
\[
\boxed{D_t=e^{-r\tau}\Q(S_T>K\mid\cF_t)
=e^{-r\tau}\Phi(d_2).}
\]
A stock-or-nothing claim paying \(S_T\1_{\{S_T>K\}}\) has value
\[
\boxed{A_t=e^{-r\tau}\E^\Q[S_T\1_{\{S_T>K\}}\mid\cF_t]
=S_t\Phi(d_1).}
\]
The second tail weights each outcome by its stock price. Completing the
square in the normal density shifts the standardized cutoff by
\(\sigma\sqrt\tau\), which is why it uses \(d_1\) instead of \(d_2\).
Combining the two claims gives \(C_t=A_t-KD_t\). The dividend-paying
replication argument below will give the corresponding result for \(q>0\).

\subsection{Forwards and call decomposition}

\begin{boardexample}{a forward payoff}
For \(H=S_T-K\) and a non-dividend-paying stock,
\[
\begin{aligned}
V_t
&=e^{-r(T-t)}\E^\Q[S_T-K\mid\cF_t]\\
&=S_t-Ke^{-r(T-t)}.
\end{aligned}
\]
The delivery price that makes the contract worth zero is therefore
\[
\boxed{K_t=S_te^{r(T-t)}.}
\]
This is both a martingale calculation and the law-of-one-price value of
\emph{one share minus a zero-coupon payment}.
\end{boardexample}

For a call,
\[
(S_T-K)^+=S_T\1_{\{S_T>K\}}-K\1_{\{S_T>K\}},
\]
so
\[
C_t=e^{-r(T-t)}
\left(
\E^\Q[S_T\1_{\{S_T>K\}}\mid\cF_t]
-K\Q(S_T>K\mid\cF_t)
\right).
\]
The first term is stock weighted; the second is a discounted cash tail. This
explains structurally why the Black--Scholes formula contains two different
normal terms.

\begin{studentexercise}{value before formula}
Let \(H=aS_T+bB_T^{(0)}\), where \(a,b\) are constants. Price the claim without
computing any distribution.
\end{studentexercise}
\begin{shortsolution}
The claim is already replicated by \(a\) shares and \(b\) bank-account units:
\[
\boxed{V_t=aS_t+bB_t^{(0)}.}
\]
The martingale identity gives the same result by linearity.
\end{shortsolution}

% ============================================================
\section{Dividend Yield and Risk-Neutral Pricing}
% ============================================================

\subsection{Continuous dividends and total gains}

If the stock pays a continuous dividend yield \(q_t\), one share earns the
total gain
\[
\boxed{\dd G_t=\dd S_t+q_tS_t\dd t.}
\]
The ex-dividend price \(S_t\) excludes cash distributions. A holder of
\(\Delta_t\) shares therefore earns \(\Delta_t\dd G_t\), not merely
\(\Delta_t\dd S_t\). This distinction matters both for replication and for
the martingale condition.

\subsection{Replicate a dividend-paying call under the historical measure}

Assume constant \(r,q,\sigma\), with \(\sigma>0\), and let \(\mu\) be the
historical \emph{ex-dividend} price drift:
\[
\dd S_t=\mu S_t\dd t+\sigma S_t\dd W_t^{\Pp},
\qquad \dd B_t=rB_t\dd t.
\]
Seek a smooth call value \(V_t=v(t,S_t)\) with
\(v(T,s)=(s-K)^+\). It\^o's formula yields
\[
\dd v=\left(v_t+\mu S_tv_s+\tfrac12\sigma^2S_t^2v_{ss}\right)\dd t
       +\sigma S_tv_s\dd W_t^{\Pp}.
\]
A self-financing portfolio holds \(\Delta_t\) shares and invests
\(v(t,S_t)-\Delta_tS_t\) in the bank account. Its gain includes the
stock's dividend cash flow:
\[
\boxed{
\dd V_t=\Delta_t(\dd S_t+qS_t\dd t)
+r(V_t-\Delta_tS_t)\dd t.}
\]
Rebalancing the holdings is financed within the portfolio; no external
cash is injected. Matching Brownian coefficients gives
\(\Delta_t=v_s(t,S_t)\). Matching the drifts then gives
\[
 v_t+\mu s v_s+\tfrac12\sigma^2s^2v_{ss}
 =\mu s v_s+qs v_s+r(v-sv_s).
\]
The historical drift cancels. The dividend-adjusted Black--Scholes
replication PDE is
\[
\boxed{
 v_t+\tfrac12\sigma^2s^2v_{ss}+(r-q)sv_s-rv=0,
 \qquad v(T,s)=(s-K)^+.}
\]
A sufficiently regular solution with an admissible self-financing hedge
replicates the call. No-arbitrage therefore identifies its price with
\(v(t,S_t)\).

\subsection{Prepaid shares, the pricing measure, and the call formula}

For constant \(q\), a prepaid forward delivering one share at \(T\)
without interim dividends is worth \(S_te^{-q(T-t)}\). Buy
\(e^{-q(T-t)}\) shares now and reinvest the dividends: the holding reaches
one share at \(T\). Consequently,
\[
\boxed{
F_{t,T}=S_te^{(r-q)(T-t)},\qquad
C_t-P_t=S_te^{-q(T-t)}-Ke^{-r(T-t)}.}
\]

The measure associated with the money-market numeraire has Brownian shift
\[
\lambda=\frac{\mu+q-r}{\sigma},\qquad
Z_t=\exp(-\lambda W_t^{\Pp}-\tfrac12\lambda^2t),\qquad
\frac{\dd\Q}{\dd\Pp}\bigg|_{\cF_T}=Z_T.
\]
Then \(W_t^{\Q}=W_t^{\Pp}+\lambda t\) is Brownian under \(\Q\), and
\[
\boxed{\frac{\dd S_t}{S_t}=(r-q)\dd t+\sigma\dd W_t^{\Q}.}
\]
In particular,
\[
\dd(D_tS_t)+D_tqS_t\dd t=D_t\sigma S_t\dd W_t^{\Q},
\qquad D_t=e^{-rt}.
\]
Thus discounted \emph{total gains} are driftless, while the discounted
ex-dividend price alone need not be a martingale.

The Feynman--Kac representation of the replication PDE is
\[
 v(t,s)=e^{-r\tau}\E^{\Q}[(S_T-K)^+\mid S_t=s],
 \qquad \tau=T-t.
\]
Conditional on \(S_t=s\), write
\[
S_T=s\exp\!\left((r-q-\tfrac12\sigma^2)\tau
             +\sigma\sqrt\tau Z\right),
\qquad Z\sim\Nlaw(0,1).
\]
Define
\[
d_2^q=\frac{\log(s/K)+(r-q-\sigma^2/2)\tau}{\sigma\sqrt\tau},
\qquad d_1^q=d_2^q+\sigma\sqrt\tau.
\]
The Gaussian cash and stock tails are
\[
\Q(S_T>K\mid S_t=s)=\Phi(d_2^q),
\qquad
 e^{-r\tau}\E^\Q[S_T\1_{\{S_T>K\}}\mid S_t=s]
 =se^{-q\tau}\Phi(d_1^q).
\]
Using \((S_T-K)^+=S_T\1_{\{S_T>K\}}-K\1_{\{S_T>K\}}\) yields the
replicating value
\[
\boxed{C_t=S_te^{-q\tau}\Phi(d_1^q)-Ke^{-r\tau}\Phi(d_2^q).}
\]
The stock holding is the derivative of this solution:
\[
\boxed{\Delta_t=v_s(t,S_t)=e^{-q\tau}\Phi(d_1^q),}
\qquad
C_t-\Delta_tS_t=-Ke^{-r\tau}\Phi(d_2^q).
\]
When differentiating, the terms involving the normal density cancel because
\(s e^{-q\tau}\varphi(d_1^q)=K e^{-r\tau}\varphi(d_2^q)\).
This price solves the PDE with the call terminal payoff, and the displayed
holdings implement the replication.

\begin{studentexercise}{spot the wrong martingale}
Under \(\Q\), suppose
\(\dd S_t/S_t=(r-q)\dd t+\sigma\dd W_t^\Q\) with \(q>0\).
Is \(e^{-rt}S_t\) a martingale? Identify the correct martingale object.
\end{studentexercise}
\begin{shortsolution}
\[
\frac{\dd(e^{-rt}S_t)}{e^{-rt}S_t}
=-q\dd t+\sigma\dd W_t^\Q,
\]
so the discounted ex-dividend price is not a martingale. The discounted
cum-dividend gain is driftless. With constant \(q\), the reinvested process
\(e^{qt}S_t\), after discounting, is a local martingale.
\end{shortsolution}

% ============================================================
\section{Changing the Numeraire}
% ============================================================

\subsection{A general numeraire}

Let \(N\) be a strictly positive traded numeraire. Under its associated
measure \(\Q^N\), prices expressed in units of \(N\) are martingales. A
replicable payoff therefore satisfies
\[
\boxed{
V_t=N_t\E^{\Q^N}\left[\frac{H}{N_T}\middle|\cF_t\right].}
\]
Changing numeraire changes both units and measure. It does not change the
economic value \(V_t\).

\begin{boardexample}{choose the numeraire that is already in the payoff}
For \(H=cN_T\),
\[
V_t=N_t\E^{\Q^N}[c\mid\cF_t]=cN_t.
\]
The best numeraire can turn a random payoff into a constant ratio.
\end{boardexample}

\begin{warningbox}{Three formulas that should not be mixed}
\[
\begin{array}{ll}
\text{money-market numeraire:}
&V_t=B_t^{(0)}\E^\Q[H/B_T^{(0)}\mid\cF_t],\\[0.25em]
\text{constant rate:}
&V_t=e^{-r(T-t)}\E^\Q[H\mid\cF_t],\\[0.25em]
\text{general numeraire:}
&V_t=N_t\E^{\Q^N}[H/N_T\mid\cF_t].
\end{array}
\]
The numeraire, measure, and discount factor are one matched package.
\end{warningbox}

% ============================================================
\section{From Prices to Hedges: Martingale Representation}
% ============================================================

\subsection{The representation theorem}

The conditional expectation gives a price process, but it does not yet state
which assets to hold. Brownian martingale representation provides the missing
link.

\begin{theorem}[Martingale Representation Theorem]
Let \((\cF_t)\) be the augmented filtration generated by a
\(d\)-dimensional \(\Q\)-Brownian motion \(W^\Q\). If
\(M\) is a square-integrable \(\Q\)-martingale, then there is a predictable
process \(\psi_t\in\R^d\) such that
\[
\boxed{
M_t=M_0+\int_0^t\psi_u^\top\dd W_u^\Q,
\qquad
\E^\Q\int_0^T\|\psi_u\|^2\dd u<\infty.}
\]
The integrand is unique up to
\(\dd t\otimes\dd\Q\)-null sets.
\end{theorem}

For a claim \(H\), take
\[
M_t=\E^\Q[D_TH\mid\cF_t]=D_tV_t.
\]
The coefficient \(\psi_t\) records the exposure of the discounted claim to
each Brownian shock.

\begin{boardexample}{linear and quadratic terminal shocks}
For \(H=W_T^\Q\),
\[
\E^\Q[H\mid\cF_t]=W_t^\Q
=\int_0^t1\dd W_u^\Q.
\]
For \(H=(W_T^\Q)^2\),
\[
\E^\Q[H\mid\cF_t]=(W_t^\Q)^2+(T-t),
\]
and It\^o's formula gives
\[
\dd\bigl((W_t^\Q)^2+T-t\bigr)=2W_t^\Q\dd W_t^\Q.
\]
Thus the representation coefficient is \(1\) in the first example and
\(2W_t^\Q\) in the second.
\end{boardexample}

\begin{studentexercise}{represent a cubic terminal shock}
Find \(M_t=\E^\Q[(W_T^\Q)^3\mid\cF_t]\) and its representation coefficient.
\end{studentexercise}
\begin{shortsolution}
Write \(W_T^\Q=W_t^\Q+Z\), where
\(Z\sim\Nlaw(0,T-t)\) is independent of \(\cF_t\). Then
\[
M_t=(W_t^\Q)^3+3(T-t)W_t^\Q.
\]
Applying It\^o's formula,
\[
\boxed{
\dd M_t=3\bigl((W_t^\Q)^2+T-t\bigr)\dd W_t^\Q.}
\]
The coefficient contains both the risk already observed and the variance
still to come.
\end{shortsolution}

\subsection{One traded asset}

Suppose the discounted stock satisfies
\[
\dd\bar S_t=\gamma_t\dd W_t^\Q,
\qquad \gamma_t\neq0,
\]
and the discounted claim martingale satisfies
\[
\dd\bar V_t=\psi_t\dd W_t^\Q.
\]
A self-financing hedge holding \(\Delta_t\) shares obeys
\[
\dd\bar V_t=\Delta_t\dd\bar S_t
=\Delta_t\gamma_t\dd W_t^\Q.
\]
Matching Brownian coefficients gives
\[
\boxed{\Delta_t=\frac{\psi_t}{\gamma_t}.}
\]
In the Black--Scholes model,
\(\gamma_t=\sigma_t\bar S_t\), so
\[
\boxed{\Delta_t=\frac{\psi_t}{\sigma_t\bar S_t}.}
\]

If \(V_t=v(t,S_t)\) is smooth and Markov, It\^o's formula yields
\[
\psi_t=D_t\sigma_tS_tv_s(t,S_t),
\]
and therefore
\[
\boxed{\Delta_t=v_s(t,S_t).}
\]
The PDE delta and the martingale-representation hedge are the same object
written in state coordinates and Brownian coordinates, respectively.

% ============================================================
\section{Two Brownian Risks and Market Completeness}
% ============================================================

\subsection{A call with an untraded sector state}

Under the historical measure, interpret \(W^1\) as an AI-sector shock and
\(W^2\) as a broad market shock. For classroom calculations suppose they
are independent. This is a simplifying modeling assumption: empirical
sector and market shocks may be correlated. Ask whether independence is
plausible and how correlation would affect the two risk directions.

Let the traded AI stock \(S\) and an untraded sector state \(Y\) satisfy
\[
\frac{\dd S_t}{S_t}
=\mu(Y_t)\dd t+\sigma_1(Y_t)\dd W_t^1+\sigma_2\dd W_t^2,
\qquad
\dd Y_t=a(Y_t)\dd t+b(Y_t)\dd W_t^1,
\]
with \(\sigma_2\neq0\). The sector state changes the stock's volatility.
For a vanilla call \(H=(S_T-K)^+\), its value is generally
\(v(t,S_t,Y_t)\). It\^o's formula gives Brownian exposure
\[
\bigl(S\sigma_1v_s+bv_y,\;S\sigma_2v_s\bigr).
\]
Holding \(\Delta\) shares of \(S\) produces only
\(\bigl(\Delta S\sigma_1,\Delta S\sigma_2\bigr)\).
The common-market component forces \(\Delta=v_s\), leaving
\(bv_y\) of sector exposure uncovered whenever it is nonzero.

\begin{warningbox}{Why a second state is needed}
If the stock coefficients were constant and the call value depended only
on \(S_t\), then \(v_y=0\). Delta hedging that call would still work even
though the stock had two Brownian terms. The example exposes
incompleteness through the untraded sector state, not simply by counting
two noises in the stock equation.
\end{warningbox}

\subsection{A second traded product and the rank condition}

Suppose a second traded asset \(R\) has dynamics under \(\Pp\)
\[
\frac{\dd R_t}{R_t}
=\nu(Y_t)\dd t+\eta_1(Y_t)\dd W_t^1+\eta_2(Y_t)\dd W_t^2.
\]
The holdings \((\Delta^S,\Delta^R)\) match the call's two shocks when
\[
\begin{pmatrix}S\sigma_1&R\eta_1\\S\sigma_2&R\eta_2\end{pmatrix}
\begin{pmatrix}\Delta^S\\\Delta^R\end{pmatrix}
=\begin{pmatrix}S\sigma_1v_s+bv_y\\S\sigma_2v_s\end{pmatrix}.
\]
If \(\sigma_1\eta_2-\sigma_2\eta_1\neq0\), the exposure matrix has full
rank and this system has a unique solution. Under the usual
integrability and admissibility conditions, full rank throughout the
modeled Brownian filtration is the spanning mechanism behind market
completeness: every suitable terminal claim can be replicated.

\subsection{The matching martingale measure}

The market-price-of-risk vector is determined by both traded assets:
\[
\begin{pmatrix}\sigma_1&\sigma_2\\\eta_1&\eta_2\end{pmatrix}
\begin{pmatrix}\lambda_1\\\lambda_2\end{pmatrix}
=\begin{pmatrix}\mu-r\\\nu-r\end{pmatrix}.
\]
Write \(W^\Pp=(W^{1,\Pp},W^{2,\Pp})^\top\) and
\(\lambda=(\lambda_1,\lambda_2)^\top\). The density process and the
terminal change of measure are
\[
\boxed{Z_t=\exp\!\left(-\int_0^t\lambda_u^\top\dd W_u^\Pp
-\frac12\int_0^t\|\lambda_u\|^2\dd u\right),\qquad
\frac{\dd\Q}{\dd\Pp}\bigg|_{\cF_T}=Z_T.}
\]
When \(Z\) is a true martingale, Girsanov gives
\(\dd W_t^{i,\Q}=\dd W_t^{i,\Pp}+\lambda_i(t)\dd t\).
Both discounted traded asset prices are then \(\Q\)-local martingales.
For the call, subject to the valuation martingale conditions,
\[
V_t=e^{-r(T-t)}\E^\Q[(S_T-K)^+\mid\cF_t].
\]
The martingale representation of the discounted claim gives two Brownian
coefficients. Matching them with the two traded exposures recovers the
same hedge as replication under \(\Pp\).

\subsection{Several assets: hedging is a linear system}

Suppose \(m\) discounted risky assets are driven by \(d\) Brownian factors:
\[
\dd\bar S_t=\Gamma_t\dd W_t^\Q,
\qquad
\Gamma_t\in\R^{m\times d}.
\]
If the claim has representation
\[
\dd\bar V_t=\psi_t^\top\dd W_t^\Q,
\qquad \psi_t\in\R^d,
\]
then holdings \(\phi_t\in\R^m\) replicate its Brownian exposure precisely
when
\[
\boxed{\Gamma_t^\top\phi_t=\psi_t.}
\]

\begin{studentexercise}{solve a two-factor hedge}
Let
\[
\Gamma=
\begin{pmatrix}
2&0\\
1&3
\end{pmatrix},
\qquad
\psi=
\begin{pmatrix}
4\\6
\end{pmatrix}.
\]
Find the asset holdings \(\phi=(\phi_1,\phi_2)^\top\).
\end{studentexercise}
\begin{shortsolution}
The hedge equation is
\[
\begin{pmatrix}2&1\\0&3\end{pmatrix}
\begin{pmatrix}\phi_1\\\phi_2\end{pmatrix}
=
\begin{pmatrix}4\\6\end{pmatrix}.
\]
Hence \(\phi_2=2\) and \(\phi_1=1\). Both Brownian shocks are spanned.
\end{shortsolution}

\begin{keyidea}{Representation versus implementation}
The MRT says every square-integrable Brownian martingale has a Brownian
exposure \(\psi\). Market completeness asks whether traded assets can
implement that exposure by solving \(\Gamma^\top\phi=\psi\). These are
different statements.
\end{keyidea}

% ============================================================
\subsection{Completeness is a spanning condition}

\begin{definition}[Complete market]
A market is complete on \([0,T]\) if every claim in the chosen admissible
class can be replicated by a self-financing trading strategy.
\end{definition}

In a Brownian diffusion market, completeness is essentially a spanning
condition. If the filtration has \(d\) independent Brownian risks and
\[
\rank(\Gamma_t)=d
\quad
\dd t\otimes\dd\Q\text{-almost everywhere},
\]
then every Brownian exposure \(\psi_t\in\R^d\) can be matched. Full column
rank requires at least \(d\) sufficiently independent risky assets.

The two-asset example above has full rank. Removing a spanning asset
leaves a Brownian risk that traded prices do not constrain.

\subsection{Incomplete markets and several martingale measures}

Consider under \(\Pp\)
\[
\frac{\dd S_t}{S_t}
=\mu\dd t+\sigma\dd W_t^1,
\]
while the filtration is generated by independent Brownian motions
\((W^1,W^2)\). To make the discounted stock driftless, the first Brownian
shift is fixed:
\[
\lambda_1=\frac{\mu-r}{\sigma}.
\]
The shift of \(W^2\) is not constrained by any traded asset. Different
choices therefore produce different equivalent martingale measures.
For a constant free shift \(\eta\), one such density is
\[
Z_T^\eta=\exp\!\left(-\lambda_1W_T^1-\eta W_T^2
-\tfrac12(\lambda_1^2+\eta^2)T\right),\qquad
\frac{\dd\Q^\eta}{\dd\Pp}\bigg|_{\cF_T}=Z_T^\eta.
\]

\begin{boardexample}{a nontraded claim reveals price ambiguity}
Let \(H=(W_T^2)^2\). Under a pricing measure indexed by a constant
\(\eta\), write
\[
W_t^{2,\Q^\eta}=W_t^2+\eta t.
\]
Then \(W_T^2=W_T^{2,\Q^\eta}-\eta T\), and
\[
\E^{\Q^\eta}[H]
=T+\eta^2T^2.
\]
The stock does not constrain \(\eta\), so martingale measures disagree on
the expectation of this nonreplicable claim. No-arbitrage alone does not
select one price.
\end{boardexample}

\begin{warningbox}{Several distinctions matter}
\begin{itemize}
  \item Multiple EMMs signal incompleteness; they are not themselves an
        arbitrage.
  \item A particular claim can have the same expectation under every EMM
        even when the market is incomplete.
  \item Redundant assets can make portfolio holdings nonunique even when
        attainable payoffs have unique prices.
  \item Completeness is about spanning payoffs, not about counting portfolio
        labels.
\end{itemize}
\end{warningbox}

\subsection{Test the rank criterion}

\begin{studentexercise}{rank, hedge, and ambiguity}
Classify each case for two Brownian factors:
\[
\Gamma_1=
\begin{pmatrix}1&0\\0&2\end{pmatrix},
\qquad
\Gamma_2=
\begin{pmatrix}1&1\\2&2\end{pmatrix}.
\]
Which matrix spans every exposure? For \(\psi=(1,-1)^\top\), does a hedge
exist in each case?
\end{studentexercise}
\begin{shortsolution}
\(\Gamma_1\) has rank two, so every exposure is spanned; the hedge solves
\(\Gamma_1^\top\phi=\psi\), giving
\(\phi=(1,-1/2)^\top\). The matrix \(\Gamma_2\) has rank one. The range of
\(\Gamma_2^\top\) is generated by \((1,1)^\top\), while
\((1,-1)^\top\) is not in that space. Equivalently,
\(\Gamma_2^\top\phi=(\phi_1+2\phi_2,\phi_1+2\phi_2)^\top\), so the second
market cannot hedge this exposure.
\end{shortsolution}

% ============================================================
\section{Asian Options and State Augmentation}
% ============================================================

We now return to the complete one-factor Black--Scholes model. Completeness
settles price uniqueness; path dependence raises a different question about
which state variables are needed to evaluate that price.

\subsection{Markov structure is optional}

If \(H=g(S_T)\) and \(S\) is Markov under \(\Q\), the conditional expectation
can often be written as \(V_t=v(t,S_t)\). This makes PDE methods available.
The pricing identity itself is more general:
\[
V_t=B_t^{(0)}
\E^\Q\left[\frac{H}{B_T^{(0)}}\middle|\cF_t\right]
\]
still makes sense when \(H\) depends on the entire path or when the state is
high-dimensional. Markov structure simplifies the representation of the
conditional expectation; it is not the source of the pricing principle.

\subsection{Path dependence changes the state, not the pricing principle}

A fixed-strike arithmetic-average call has payoff
\[
H=\left(\frac1T\int_0^TS_u\dd u-K\right)^+.
\]
Two stock paths can have the same endpoint and different averages, so
\(S_t\) alone does not contain enough information to value the claim.
Nevertheless,
\[
\boxed{
V_t=B_t^{(0)}
\E^\Q\left[\frac{H}{B_T^{(0)}}\middle|\cF_t\right]}
\]
remains valid. Path dependence complicates the state and computation, not
the martingale valuation identity.

\begin{studentexercise}{same endpoint, different claim value}
On \([0,1]\), compare the deterministic paths
\[
S_u^{(1)}=100,
\qquad
S_u^{(2)}=
\begin{cases}
140,&u<3/4,\\
60,&3/4\leq u<1,\\
100,&u=1.
\end{cases}
\]
They have the same endpoint. Compare their arithmetic averages and explain
why a terminal-state function \(v(t,S_t)\) cannot distinguish them.
\end{studentexercise}
\begin{shortsolution}
The first average is \(100\). Ignoring the single endpoint, which has zero
Lebesgue measure, the second average is
\[
\frac34(140)+\frac14(60)=120.
\]
For a strike between \(100\) and \(120\), the two paths generate different
Asian payoffs despite having the same terminal stock value. Endpoint
information alone cannot determine the running integral.
\end{shortsolution}

\subsection{The running integral restores Markov structure}

Introduce
\[
\boxed{A_t=\int_0^tS_u\dd u,\qquad \dd A_t=S_t\dd t.}
\]
Then
\[
H=\left(\frac{A_T}{T}-K\right)^+,
\]
and the enlarged state is \((S_t,A_t)\). Under constant \(r,q,\sigma\),
\[
\frac{\dd S_t}{S_t}=(r-q)\dd t+\sigma\dd W_t^\Q,
\qquad
\dd A_t=S_t\dd t.
\]
The price may therefore be written
\[
v(t,s,a)
=e^{-r(T-t)}
\E^\Q\left[
\left(\frac{A_T}{T}-K\right)^+
\middle|S_t=s,\ A_t=a
\right].
\]

\begin{proposition}[Two-state Asian PDE]
Under the preceding assumptions, a sufficiently smooth value function solves
\[
\boxed{
v_t+(r-q)sv_s+sv_a
+\frac12\sigma^2s^2v_{ss}-rv=0,}
\]
with terminal condition
\[
\boxed{v(T,s,a)=\left(\frac aT-K\right)^+.}
\]
\end{proposition}

\begin{proof}
Apply the two-dimensional It\^o formula:
\[
\dd v
=\left(v_t+(r-q)sv_s+sv_a
+\frac12\sigma^2s^2v_{ss}\right)\dd t
+\sigma sv_s\dd W_t^\Q.
\]
Discounting and setting the drift to zero gives the PDE. There is no
\(v_{aa}\) term because \(A\) has finite variation and no direct Brownian
shock.
\end{proof}

\begin{warningbox}{Do not invent diffusion terms}
The state has two coordinates but only one Brownian driver. Since
\(\dd A_t=S_t\dd t\),
\[
\dd[A]_t=0,\qquad \dd[S,A]_t=0.
\]
State dimension and noise dimension are not the same thing.
\end{warningbox}

\subsection{Arithmetic, geometric, and discrete monitoring}

For a geometric-average option,
\[
G_T=\exp\left(\frac1T\int_0^T\log S_u\dd u\right).
\]
In Black--Scholes, \(\log G_T\) is Gaussian because it is a linear functional
of Brownian motion. This often yields a closed form. The arithmetic average
is an integral of a lognormal process and generally has no elementary
distribution.

For monitoring dates \(t_1<\cdots<t_n\), define
\[
A_k=\sum_{j=1}^kS_{t_j}.
\]
Then
\[
A_{k+1}=A_k+S_{t_{k+1}},
\qquad
H=\left(\frac{A_n}{n}-K\right)^+.
\]
Between fixing dates, \(A_k\) is constant; at a fixing date, it jumps by the
new observation. Dynamic programming therefore uses the state
\((S_{t_k},A_k,k)\).

\begin{boardexample}{Monte Carlo workflow for an arithmetic Asian call}
For \(M\) simulated paths:
\begin{enumerate}
  \item simulate the stock under \(\Q\);
  \item update the running sum or time integral along each path;
  \item compute \(H^{(m)}\) at maturity;
  \item estimate
  \[
  \widehat V_0=e^{-rT}\frac1M\sum_{m=1}^MH^{(m)};
  \]
  \item report a standard error
  \[
  e^{-rT}\frac{\widehat{\operatorname{sd}}(H)}{\sqrt M}.
  \]
\end{enumerate}
Path dependence is computationally natural for simulation because the full
path is generated anyway.
\end{boardexample}

\begin{studentexercise}{select the minimal state}
For each payoff, propose a Markov state augmentation:
\[
\max_{0\leq u\leq T}S_u-K,\qquad
\1_{\{\min_{0\leq u\leq T}S_u>L\}}(S_T-K)^+,\qquad
\left(\frac1T\int_0^TS_u\dd u-K\right)^+.
\]
\end{studentexercise}
\begin{shortsolution}
Track respectively the running maximum
\(M_t=\max_{u\leq t}S_u\), the running minimum
\(m_t=\min_{u\leq t}S_u\) or a barrier-survival indicator, and the running
integral \(A_t=\int_0^tS_u\dd u\). The state should summarize precisely the
past information that the remaining payoff needs.
\end{shortsolution}

% ============================================================
\section{Synthesis: Choosing the Right Route}
% ============================================================

\subsection{Choose the right concept before calculating}

\begin{studentexercise}{diagnose six situations}
Match each situation with the main concept it tests:
\begin{enumerate}
  \item discounted traded prices retain a nonzero drift under a candidate
        measure;
  \item two portfolios have identical terminal cashflows but different costs;
  \item a claim's Brownian exposure is outside the span of \(\Gamma_t^\top\);
  \item exercise is allowed before maturity;
  \item a payoff depends on \(\int_0^TS_u\dd u\);
  \item several equivalent martingale measures price a nontraded risk
        differently.
\end{enumerate}
\end{studentexercise}
\begin{shortsolution}
\begin{enumerate}
  \item the measure is not an EMM for the chosen numeraire;
  \item law of one price and an arbitrage opportunity;
  \item nonreplicability and market incompleteness;
  \item optimal stopping and the Snell envelope;
  \item path dependence and state augmentation;
  \item incomplete-market price ambiguity.
\end{enumerate}
\end{shortsolution}

\subsection{A complete martingale-pricing drill}

\begin{studentexercise}{from payoff to hedge}
Suppose \(r\) is constant and
\[
\dd\bar S_t=\sigma\bar S_t\dd W_t^\Q,\qquad \sigma>0.
\]
Let \(H\in L^2(\Q)\) be replicable and define
\[
M_t=\E^\Q[e^{-rT}H\mid\cF_t].
\]
Assume the MRT representation is
\[
\dd M_t=\psi_t\dd W_t^\Q.
\]
\begin{enumerate}
  \item Write the claim price \(V_t\).
  \item Find the number of stock shares in the hedge.
  \item State the extra property that makes the hedge implementation valid.
\end{enumerate}
\end{studentexercise}
\begin{shortsolution}
\[
\boxed{
V_t=e^{rt}M_t
=e^{-r(T-t)}\E^\Q[H\mid\cF_t].}
\]
Since \(M_t=e^{-rt}V_t=\bar V_t\),
\[
\dd\bar V_t=\psi_t\dd W_t^\Q
=\Delta_t\sigma\bar S_t\dd W_t^\Q,
\]
so
\[
\boxed{\Delta_t=\frac{\psi_t}{\sigma\bar S_t}.}
\]
The Brownian exposure must be spanned by traded assets, and the strategy must
be predictable, self-financing, admissible, and sufficiently integrable.
\end{shortsolution}

\subsection{Pricing checklist}

\begin{keyidea}{Eight checks for every martingale-pricing problem}
\begin{enumerate}
  \item What is the numeraire and therefore the discount factor?
  \item Under which measure are discounted traded gains martingales?
  \item Is the relevant local martingale a true martingale?
  \item Is the payoff replicable, or is some risk unspanned?
  \item Which Brownian exposures appear in the MRT coefficient?
  \item Do the traded assets span those exposures?
  \item Which state variables summarize the remaining payoff?
  \item Does exercise introduce an optimization over stopping times?
\end{enumerate}
\end{keyidea}

\begin{keyidea}{Final synthesis}
\begin{itemize}
  \item Discounted self-financing wealth is a stochastic integral in
        discounted traded gains.
  \item Conditional expectation computes the value of a replicable claim
        once a numeraire and its martingale measure are fixed.
  \item The MRT reveals Brownian hedge exposures; the asset volatility matrix
        determines whether those exposures can be traded.
  \item Completeness is a spanning property. Unspanned risks may leave
        several martingale measures and ambiguous values for nontraded claims.
  \item American options replace one terminal expectation by optimal stopping
        and a Snell envelope.
  \item Asian options keep the same pricing identity but require an enlarged
        state or path simulation.
\end{itemize}
\end{keyidea}

% ============================================================
\section{Revision Exercises: Lectures 2--6}

Four exercises revisit tools needed throughout the course. Each question
appears on a single slide with its short answer; the calculations below
give the intermediate steps.

\subsection{It\^o's formula for a nonlinear function}

\begin{studentexercise}{the sine of an It\^o process}
Let \(X_0=0\) and
\[
\dd X_t=(1+X_t)\dd t+(1+t)\dd W_t.
\]
Find the stochastic differential of \(Y_t=\sin(X_t)\), identifying the
quadratic-variation correction.
\end{studentexercise}
\begin{shortsolution}
For \(f(x)=\sin x\), we have \(f'(x)=\cos x\) and
\(f''(x)=-\sin x\). The finite-variation drift contributes no quadratic
variation, while
\[
\dd[X]_t=(1+t)^2\dd t.
\]
It\^o's formula therefore gives
\[
\begin{aligned}
\dd\sin(X_t)
&=f'(X_t)\dd X_t+\tfrac12f''(X_t)\dd[X]_t\\
&=\boxed{\left[(1+X_t)\cos X_t
 -\tfrac12(1+t)^2\sin X_t\right]\dd t
 +(1+t)\cos X_t\dd W_t.}
\end{aligned}
\]
The negative sine term is created by curvature and quadratic variation; it
would be absent under the ordinary chain rule.
\end{shortsolution}

\subsection{Girsanov for a deterministic time-varying drift}

\begin{studentexercise}{make the drifted process Brownian}
Under \(\Pp\), let \(Y_0=0\) and
\[
\dd Y_t=(\alpha+\beta t)\dd t+\dd W_t,
\qquad 0\leq t\leq T.
\]
Construct a density defining an equivalent measure \(\Q\) under which
\(Y\) is Brownian. Explain why the density is a true martingale.
\end{studentexercise}
\begin{shortsolution}
Set \(\theta_t=\alpha+\beta t\) and define
\[
Z_t=\exp\!\left(-\int_0^t\theta_s\dd W_s
 -\tfrac12\int_0^t\theta_s^2\dd s\right),
\qquad
\frac{\dd\Q}{\dd\Pp}\bigg|_{\cF_T}=Z_T.
\]
The deterministic drift is square-integrable on every finite horizon:
\[
\int_0^T\theta_s^2\dd s
=\alpha^2T+\alpha\beta T^2+\frac{\beta^2T^3}{3}<\infty.
\]
Thus Novikov's condition holds and \(Z\) is a true martingale. By
Girsanov,
\[
W_t^\Q=W_t+\int_0^t\theta_s\dd s
=W_t+\alpha t+\tfrac12\beta t^2
\]
is Brownian under \(\Q\). Since \(Y_0=0\), integration of the original
SDE gives \(Y_t=W_t^\Q\), equivalently \(\dd Y_t=\dd W_t^\Q\).
\end{shortsolution}

\subsection{The maximum of Brownian motion with drift}

\begin{studentexercise}{Girsanov plus reflection}
Let
\[
M_T^\mu=\max_{0\leq s\leq T}(W_s+\mu s),
\qquad T>0,\quad a>0.
\]
Use a change of measure and the reflection principle to compute
\(\Pp(M_T^\mu\geq a)\). Check the result when \(\mu=0\).
\end{studentexercise}
\begin{shortsolution}
Set \(X_t=W_t+\mu t\) and
\[
Z_T=\exp(-\mu W_T-\tfrac12\mu^2T),
\qquad
\frac{\dd\Q}{\dd\Pp}\bigg|_{\cF_T}=Z_T.
\]
Under \(\Q\), \(X\) is a standard Brownian motion. Rewriting the inverse
density in terms of \(X_T\) gives
\[
\frac{\dd\Pp}{\dd\Q}\bigg|_{\cF_T}
=\exp(\mu X_T-\tfrac12\mu^2T).
\]
Let \(\varphi_T(x)=(2\pi T)^{-1/2}e^{-x^2/(2T)}\) and
\(M_T=\max_{s\leq T}X_s\) under \(\Q\). For \(x<a\), reflection at the first
hit of \(a\) gives
\[
\Q(M_T\geq a,\ X_T\in\dd x)=\varphi_T(2a-x)\dd x.
\]
For \(x\geq a\), the event \(X_T\geq a\) already implies a hit. Therefore
\[
\begin{aligned}
\Pp(M_T^\mu\geq a)
=e^{-\mu^2T/2}\bigg[
&\int_a^\infty e^{\mu x}\varphi_T(x)\dd x\\
+&\int_{-\infty}^a e^{\mu x}\varphi_T(2a-x)\dd x
\bigg].
\end{aligned}
\]
Completing the square in the first integral and substituting \(y=2a-x\)
in the second gives
\[
\boxed{
\Pp(M_T^\mu\geq a)
=\Phi\!\left(\frac{\mu T-a}{\sqrt T}\right)
+e^{2\mu a}\Phi\!\left(\frac{-\mu T-a}{\sqrt T}\right).}
\]
Here \(\Phi\) is the standard normal CDF. When \(\mu=0\), the formula
reduces to \(2[1-\Phi(a/\sqrt T)]\), as required by ordinary reflection.
\end{shortsolution}

\subsection{Stopping times versus times that use future information}

\begin{studentexercise}{classify five random times}
Use the Brownian filtration \((\cF_t)\). Fix \(T>0\) and \(a>0\). Decide
which of the following are stopping times:
\[
\begin{array}{ll}
\tau_a=\inf\{t\geq0:W_t\geq a\}, & T/2,\\[0.3em]
\eta=\inf\{t\geq T/2:W_t=0\}, &
\rho_T=\sup\{t\leq T:W_t=0\},\\[0.3em]
\gamma_T=\operatorname*{arg\,max}_{0\leq s\leq T}W_s.&
\end{array}
\]
\end{studentexercise}
\begin{shortsolution}
The first three are stopping times:
\begin{itemize}
  \item \(\{\tau_a\leq t\}
  =\{\max_{0\leq s\leq t}W_s\geq a\}\in\cF_t\);
  \item \(T/2\) is deterministic;
  \item for \(t\geq T/2\), \(\{\eta\leq t\}\) says that the observed path
        has a zero on \([T/2,t]\); for \(t<T/2\), the event is empty.
\end{itemize}
The last zero \(\rho_T\) and the time of the maximum \(\gamma_T\) are not
stopping times: at a candidate time \(t<T\), one must inspect the remaining
path to rule out a later zero or a higher value. The criterion throughout is
\(\{\tau\leq t\}\in\cF_t\) for every \(t\).
\end{shortsolution}

\end{document}
